\subsection{Approaching from a side and growing without bound}

\begin{definitionbox}[One-sided limits]
The symbols
\[
\lim_{x\to a^-}f(x)
\qquad\text{and}\qquad
\lim_{x\to a^+}f(x)
\]
describe approach from values smaller and larger than $a$, respectively.
\end{definitionbox}

A two-sided finite limit exists only when both one-sided limits exist and are equal.

\begin{examplebox}[A jump]
For
\[
f(x)=\begin{cases}0,&x<0,\\1,&x\ge0,\end{cases}
\]
we have
\[
\lim_{x\to0^-}f(x)=0,
\qquad
\lim_{x\to0^+}f(x)=1.
\]
Therefore the two-sided limit at zero does not exist.
\end{examplebox}

\begin{examplebox}[A vertical asymptote]
For $f(x)=1/x$,
\[
\lim_{x\to0^+}\frac1x=\infty,
\qquad
\lim_{x\to0^-}\frac1x=-\infty.
\]
The line $x=0$ is a vertical asymptote.
\end{examplebox}

\begin{interpretationbox}[Infinite limits]
The symbol $\infty$ does not denote a real output value. It records unbounded growth of the function values.
\end{interpretationbox}
